% GENERATED FILE. Edit canonical lessons, then run npm run generate:books.
% canonical-source-hash: fnv1a64:c9488218df351546
% canonical-lessons: PA-C274-minna, PA-C274-gaz, PA-C274-mitar, PA-C274-nap, PA-C274-jora

\chapter{Yards and Metres}
\label{ch:pa-a2-yards-and-metres}
\hlchaptermodality{274}

\begin{chapteropening}
\textbf{By the end of this chapter:} \emph{I can say how long or wide something is.}

Complete the last lesson of chapter 274: I can say how long or wide something is.
\end{chapteropening}

\section[\texorpdfstring{\pa{ਮਿਣਨਾ} (miṇnā)}{ਮਿਣਨਾ (miṇnā)}]{\pa{ਮਿਣਨਾ} (miṇnā) — to measure}

\label{lesson:PA-C274-minna}



\begin{quote}
Before the new one: say the Punjabi for weight, then the Punjabi for to carry a load.
\end{quote}

\subsection*{You'll want to know: \pa{ਮਿਣਨਾ}}
\textbf{\pa{ਮਿਣਨਾ}} — \emph{miṇnā} — \textquotedblleft{}to measure\textquotedblright{}.

\begin{quote}
\pa{ਦਰਜ਼ੀ} \pa{ਮੇਰੀ} \pa{ਕਮੀਜ਼} \pa{ਮਿਣਦਾ} \pa{ਹੈ।}
\end{quote}

\begin{quote}
\emph{darzī merī kamīz miṇdā hai.}
\end{quote}

\begin{quote}
The tailor measures my shirt.
\end{quote}

Past (I …): \pa{ਮੈਂ} \pa{ਕਮਰਾ} \pa{ਮਿਣਿਆ}, \emph{maiṁ kamrā miṇiā}.

\begin{cousinweb}
\textbf{\pa{ਮਿਣਨਾ}} (\emph{miṇnā}) is inherited from a reconstructed \emph{*mināti}, \textquotedblleft{}measures\textquotedblright{}; CDIAL lists Punjabi \emph{miṇnā}, \textquotedblleft{}to measure\textquotedblright{}, under it (CDIAL 10132).
\end{cousinweb}

\begin{grammarlens}[title={What you've built}]
One more word for this chapter.
\end{grammarlens}

\subsection*{Your turn}
\begin{itemize}
\raggedright
  \item \emph{Say it:} \emph{miṇnā}
  \item \emph{Say it:} \emph{miṇnā}, once more
  \item \emph{Say it:} \emph{darzī merī kamīz miṇdā hai.} — The tailor measures my shirt.
  \item \emph{Recall:} say the Punjabi for weight, then the Punjabi for to carry a load, then say \emph{miṇnā} again
\end{itemize}

\begin{tcolorbox}[breakable,colback=teal!4,colframe=teal!35!black,title={Before you move on}]
What does \pa{ਮਿਣਨਾ} mean? (\textquotedblleft{}To measure\textquotedblright{}.) Say it once more.
\end{tcolorbox}
\section[\texorpdfstring{\pa{ਗਜ਼} (gaz)}{ਗਜ਼ (gaz)}]{\pa{ਗਜ਼} (gaz) — a yard}

\label{lesson:PA-C274-gaz}



\begin{quote}
Before the new one: say the Punjabi for to carry a load, then the Punjabi for to measure.
\end{quote}

\subsection*{You'll want to know: \pa{ਗਜ਼}}
\textbf{\pa{ਗਜ਼}} — \emph{gaz} — \textquotedblleft{}a yard\textquotedblright{}.

\begin{quote}
\pa{ਮੇਰਾ} \pa{ਕਮਰਾ} \pa{ਪੰਜ} \pa{ਗਜ਼} \pa{ਲੰਮਾ} \pa{ਹੈ।}
\end{quote}

\begin{quote}
\emph{merā kamrā panj gaz lammā hai.}
\end{quote}

\begin{quote}
My room is five yards long.
\end{quote}

A masculine noun.

\begin{cousinweb}
\textbf{\pa{ਗਜ਼}} (\emph{gaz}) is a loan, usually traced to Persian \emph{gaz}, \textquotedblleft{}a yard, an ell\textquotedblright{}.
\end{cousinweb}

\begin{grammarlens}[title={What you've built}]
One more word for this chapter.
\end{grammarlens}

\subsection*{Your turn}
\begin{itemize}
\raggedright
  \item \emph{Say it:} \emph{gaz}
  \item \emph{Say it:} \emph{gaz}, once more
  \item \emph{Say it:} \emph{merā kamrā panj gaz lammā hai.} — My room is five yards long.
  \item \emph{Recall:} say the Punjabi for to carry a load, then the Punjabi for to measure, then say \emph{gaz} again
  \item \emph{Write it:} copy \textbf{\pa{ਗਜ਼}} once, checking the whole word
\end{itemize}

\begin{tcolorbox}[breakable,colback=teal!4,colframe=teal!35!black,title={Before you move on}]
What does \pa{ਗਜ਼} mean? (\textquotedblleft{}A yard\textquotedblright{}.) Say it once more.
\end{tcolorbox}
\section[\texorpdfstring{\pa{ਮੀਟਰ} (mīṭar)}{ਮੀਟਰ (mīṭar)}]{\pa{ਮੀਟਰ} (mīṭar) — a metre}

\label{lesson:PA-C274-mitar}



\begin{quote}
Before the new one: say the Punjabi for to measure, then the Punjabi for a yard.
\end{quote}

\subsection*{You'll want to know: \pa{ਮੀਟਰ}}
\textbf{\pa{ਮੀਟਰ}} — \emph{mīṭar} — \textquotedblleft{}a metre\textquotedblright{}.

\begin{quote}
\pa{ਸੜਕ} \pa{ਦਸ} \pa{ਮੀਟਰ} \pa{ਚੌੜੀ} \pa{ਹੈ।}
\end{quote}

\begin{quote}
\emph{saṛak das mīṭar chauṛī hai.}
\end{quote}

\begin{quote}
The road is ten metres wide.
\end{quote}

A masculine noun.

\begin{cousinweb}
\textbf{\pa{ਮੀਟਰ}} (\emph{mīṭar}) is borrowed from English \emph{metre}.
\end{cousinweb}

\begin{grammarlens}[title={What you've built}]
One more word for this chapter.
\end{grammarlens}

\subsection*{Your turn}
\begin{itemize}
\raggedright
  \item \emph{Say it:} \emph{mīṭar}
  \item \emph{Say it:} \emph{mīṭar}, once more
  \item \emph{Say it:} \emph{saṛak das mīṭar chauṛī hai.} — The road is ten metres wide.
  \item \emph{Recall:} say the Punjabi for to measure, then the Punjabi for a yard, then say \emph{mīṭar} again
  \item \emph{Write it:} copy \textbf{\pa{ਮੀਟਰ}} once, checking the whole word
\end{itemize}

\begin{tcolorbox}[breakable,colback=teal!4,colframe=teal!35!black,title={Before you move on}]
What does \pa{ਮੀਟਰ} mean? (\textquotedblleft{}A metre\textquotedblright{}.) Say it once more.
\end{tcolorbox}
\section[\texorpdfstring{\pa{ਨਾਪ} (nāp)}{ਨਾਪ (nāp)}]{\pa{ਨਾਪ} (nāp) — a size; a measurement}

\label{lesson:PA-C274-nap}



\begin{quote}
Before the new one: say the Punjabi for a yard, then the Punjabi for a metre.
\end{quote}

\subsection*{You'll want to know: \pa{ਨਾਪ}}
\textbf{\pa{ਨਾਪ}} — \emph{nāp} — \textquotedblleft{}a size; a measurement\textquotedblright{}.

\begin{quote}
\pa{ਦਰਜ਼ੀ} \pa{ਮੇਰਾ} \pa{ਨਾਪ} \pa{ਲੈਂਦਾ} \pa{ਹੈ।}
\end{quote}

\begin{quote}
\emph{darzī merā nāp laindā hai.}
\end{quote}

\begin{quote}
The tailor takes my measurements.
\end{quote}

A masculine noun.

\begin{cousinweb}
\textbf{\pa{ਨਾਪ}} (\emph{nāp}) goes with \emph{nāpṇā}, \textquotedblleft{}to measure\textquotedblright{}; CDIAL lists Punjabi \emph{nāpṇā} under Sanskrit \emph{jñāpyate}, \textquotedblleft{}is made known\textquotedblright{} (CDIAL 5283).
\end{cousinweb}

\begin{grammarlens}[title={What you've built}]
One more word for this chapter.
\end{grammarlens}

\subsection*{Your turn}
\begin{itemize}
\raggedright
  \item \emph{Say it:} \emph{nāp}
  \item \emph{Say it:} \emph{nāp}, once more
  \item \emph{Say it:} \emph{darzī merā nāp laindā hai.} — The tailor takes my measurements.
  \item \emph{Recall:} say the Punjabi for a yard, then the Punjabi for a metre, then say \emph{nāp} again
  \item \emph{Write it:} copy \textbf{\pa{ਨਾਪ}} once, checking the whole word
\end{itemize}

\begin{tcolorbox}[breakable,colback=teal!4,colframe=teal!35!black,title={Before you move on}]
What does \pa{ਨਾਪ} mean? (\textquotedblleft{}A size; a measurement\textquotedblright{}.) Say it once more.
\end{tcolorbox}
\section[\texorpdfstring{\pa{ਜੋੜਾ} (joṛā)}{ਜੋੜਾ (joṛā)}]{\pa{ਜੋੜਾ} (joṛā) — a pair}

\label{lesson:PA-C274-jora}



\begin{quote}
Before the new one: say the Punjabi for a metre, then the Punjabi for a size.
\end{quote}

\subsection*{You'll want to know: \pa{ਜੋੜਾ}}
\textbf{\pa{ਜੋੜਾ}} — \emph{joṛā} — \textquotedblleft{}a pair\textquotedblright{}.

\begin{quote}
\pa{ਮੈਂ} \pa{ਜੁੱਤੀਆਂ} \pa{ਦਾ} \pa{ਇੱਕ} \pa{ਨਵਾਂ} \pa{ਜੋੜਾ} \pa{ਖ਼ਰੀਦਾਂਗਾ।}
\end{quote}

\begin{quote}
\emph{maiṁ juttīāṁ dā ikk navāṁ joṛā kharīdāṅgā.}
\end{quote}

\begin{quote}
I will buy a new pair of shoes.
\end{quote}

A masculine noun.

\begin{cousinweb}
\textbf{\pa{ਜੋੜਾ}} (\emph{joṛā}) is inherited from a reconstructed \emph{*yuṭati}, \textquotedblleft{}is joined\textquotedblright{}; CDIAL lists Punjabi \emph{joṛā}, \textquotedblleft{}couple\textquotedblright{}, under it (CDIAL 10496).
\end{cousinweb}

\begin{grammarlens}[title={What you've built}]
One more word for this chapter.
\end{grammarlens}

\subsection*{Your turn}
\begin{itemize}
\raggedright
  \item \emph{Say it:} \emph{joṛā}
  \item \emph{Say it:} \emph{joṛā}, once more
  \item \emph{Say it:} \emph{maiṁ juttīāṁ dā ikk navāṁ joṛā kharīdāṅgā.} — I will buy a new pair of shoes.
  \item \emph{Recall:} say the Punjabi for a metre, then the Punjabi for a size, then say \emph{joṛā} again
  \item \emph{Write it:} copy \textbf{\pa{ਜੋੜਾ}} once, checking the whole word
\end{itemize}

\begin{tcolorbox}[breakable,colback=teal!4,colframe=teal!35!black,title={Before you move on}]
What does \pa{ਜੋੜਾ} mean? (\textquotedblleft{}A pair\textquotedblright{}.) Say it once more.
\end{tcolorbox}
